\documentclass[11pt,a4paper]{article}

% --- Core Packages ---
\usepackage[utf8]{inputenc}
\usepackage[T1]{fontenc}
\usepackage{lmodern}
\usepackage{microtype}
\usepackage[margin=1in]{geometry}
\usepackage{amsmath,amssymb,amsthm}
\usepackage{booktabs}
\usepackage{tabularx}
\usepackage{array}
\usepackage{graphicx}
\usepackage{float}
\usepackage{xcolor}
\usepackage{enumitem}
\usepackage{caption}
\usepackage{fancyhdr}
\usepackage[colorlinks=true,linkcolor=blue!60!black,citecolor=blue!60!black,urlcolor=blue!60!black]{hyperref}

% --- Theorem & Proof Environments ---
\newtheorem{proposition}{Proposition}
\newtheorem{remark}{Remark}

% --- Header & Footer Styling ---
\pagestyle{fancy}
\fancyhf{}
\lhead{\small\textit{PID Controller Optimization via PSO Algorithm}}
\rhead{\small\textit{Control Systems \& Swarm Intelligence}}
\cfoot{\thepage}
\setlength{\headheight}{14pt}

% --- Helper Macro for Conditional Figure Inclusion ---
\newcommand{\includeResultFigure}[3]{%
  \begin{figure}[H]
    \centering
    \IfFileExists{#1}{%
      \includegraphics[width=0.86\textwidth]{#1}%
    }{%
      \fbox{\parbox{0.82\textwidth}{\centering\vspace{1.2cm}%
        \textbf{Figure Pending Simulation Output}\\[4pt]
        \texttt{#1}\\[4pt]
        \small Run \texttt{src/pso\_pid\_tuning.m} in MATLAB to generate this 300-DPI plot, then recompile.%
      \vspace{1.2cm}}}%
    }
    \caption{#2}
    \label{#3}
  \end{figure}%
}

\begin{document}

% --- Title Block ---
\begin{center}
  {\LARGE\bfseries PID Controller Optimization via Particle Swarm Optimization (PSO) Algorithm}\\[0.6em]
  {\large Automated Gain Synthesis, Analytical Stability Proofs, and Closed-Loop Step/Disturbance Validation in MATLAB/Simulink}\\[0.8em]
  {\normalsize\textbf{Project Timeline:} May 2025 -- November 2025 \quad|\quad \textbf{Platform:} MATLAB \& Simulink}\\[0.4em]
  \rule{\textwidth}{0.8pt}
\end{center}

\begin{abstract}
Proportional--Integral--Derivative (PID) controllers govern the majority of industrial feedback loops, yet classical tuning rules such as the Ziegler--Nichols ultimate-sensitivity method frequently yield excessive overshoot and narrow phase margins on higher-order plants. This report presents an automated, stability-constrained Particle Swarm Optimization (PSO) framework for synthesizing optimal PID parameters $(K_p, K_i, K_d)$ on a benchmark third-order linear time-invariant (LTI) system $G(s) = \frac{1}{(s+1)(s+2)(s+3)}$. Analytical proofs establish the Routh--Hurwitz marginal stability boundary ($K_u = 60$, $P_u = 2\pi/\sqrt{11}\text{ s}$), derivative-filter realizability ($T_f = 0.01\text{ s}$), and asymptotic zero steady-state error under both unit-step setpoint tracking and step load disturbances via the Final Value Theorem. Benchmark validation in MATLAB and Simulink compares the PSO-tuned controller against classical Ziegler--Nichols tuning across time-domain transient metrics, $s$-plane closed-loop pole locations, and frequency-domain gain and phase stability margins.
\end{abstract}

% ==============================================================================
\section{Introduction and Problem Formulation}
% ==============================================================================

Despite advances in modern state-space and predictive control, the three-term PID controller remains the standard industrial workhorse due to its structural simplicity and physical interpretability. However, tuning the proportional ($K_p$), integral ($K_i$), and derivative ($K_d$) gains involves fundamental engineering trade-offs: increasing $K_p$ and $K_i$ accelerates rise time and eliminates steady-state offset, but degrades relative stability, increases maximum percentage overshoot ($M_p$), and prolongs settling time ($t_s$).

Classical empirical tuning rules—most notably the closed-loop Ziegler--Nichols method—derive gains directly from the plant's ultimate gain ($K_u$) and ultimate oscillation period ($P_u$). Because Ziegler--Nichols rules are calibrated for a quarter-amplitude decay ratio, they push closed-loop poles close to the imaginary axis, resulting in pronounced oscillatory ringing ($M_p \approx 40\%\text{--}60\%$) and narrow phase margins.

To overcome these limitations, this project formulates PID gain synthesis as a constrained multi-objective optimization problem solved via Particle Swarm Optimization (PSO):
\begin{enumerate}[leftmargin=2em]
  \item \textbf{Automated Gain Synthesis:} Evolve a 35-particle swarm across a continuous 3D gain space to minimize a weighted combination of Integral of Time-weighted Absolute Error (ITAE), maximum overshoot, settling time, and steady-state error.
  \item \textbf{Guaranteed Closed-Loop Stability:} Enforce Hurwitz pole stability at every candidate evaluation via an exterior penalty barrier.
  \item \textbf{Dual-Domain Verification:} Validate transient tracking, external load disturbance rejection ($t = 4\text{ s}$), closed-loop $s$-plane pole migration, and open-loop Gain/Phase margins in MATLAB and Simulink.
\end{enumerate}

% ==============================================================================
\section{Mathematical Modeling and Analytical Stability Proofs}
% ==============================================================================

\subsection{Benchmark Third-Order Plant Dynamics}
The controlled process modeled in \texttt{src/pso\_pid\_tuning.m} and \texttt{simulink/pid\_pso\_demo.slx} is an overdamped third-order LTI plant with three distinct real open-loop poles at $s = -1, -2, -3$:
\begin{equation}
  G(s) = \frac{1}{(s+1)(s+2)(s+3)} = \frac{1}{s^3 + 6s^2 + 11s + 6}.
  \label{eq:plant}
\end{equation}
Because all open-loop poles satisfy $\operatorname{Re}(s) < 0$, $G(s)$ is bounded-input bounded-output (BIBO) stable with DC gain $G(0) = \frac{1}{6}$ and system \textbf{Type 0} (zero open-loop integrators). Under pure proportional feedback, a unit-step setpoint produces a nonzero steady-state offset, necessitating integral action.

\subsection{Routh--Hurwitz Proof of Ultimate Gain (\texorpdfstring{$K_u$}{Ku}) and Period (\texorpdfstring{$P_u$}{Pu})}
To establish the classical Ziegler--Nichols baseline, consider unity negative feedback with a pure proportional controller $C(s) = K$. The closed-loop characteristic equation $1 + K G(s) = 0$ expands to:
\begin{equation}
  s^3 + 6s^2 + 11s + (6 + K) = 0.
  \label{eq:char_eq}
\end{equation}

\begin{proposition}[Ziegler--Nichols Marginal Stability Parameters]
For the plant $G(s) = \frac{1}{(s+1)(s+2)(s+3)}$ under unity proportional feedback, closed-loop asymptotic stability holds if and only if $-6 < K < 60$. At the upper stability boundary, the ultimate gain is $K_u = 60$ and the ultimate oscillation period is $P_u = \frac{2\pi}{\sqrt{11}}\text{ s} \approx 1.8944\text{ s}$.
\end{proposition}

\begin{proof}
Constructing the Routh--Hurwitz array for the cubic polynomial \eqref{eq:char_eq}:
\[
\begin{array}{c|cc}
s^3 & 1 & 11 \\
s^2 & 6 & 6 + K \\
s^1 & b_1 & 0 \\
s^0 & c_1 & 0
\end{array}
\]
where the first-column elements are:
\[
  b_1 = \frac{(6)(11) - (1)(6 + K)}{6} = \frac{60 - K}{6}, \qquad c_1 = 6 + K.
\]
By the Routh--Hurwitz stability criterion, all roots lie in the open left-half $s$-plane if and only if every element in the first column is strictly positive:
\[
  6 + K > 0 \implies K > -6, \qquad \frac{60 - K}{6} > 0 \implies K < 60.
\]
Setting $b_1 = 0$ yields the ultimate gain $K_u = 60$ at which a complex-conjugate pole pair crosses the imaginary axis. Forming the auxiliary polynomial $A(s)$ from the $s^2$ row at $K = 60$:
\[
  A(s) = 6s^2 + (6 + 60) = 6s^2 + 66 = 0 \implies s^2 = -11 \implies s = \pm j\sqrt{11}.
\]
Thus, the sustained oscillation frequency is $\omega_u = \sqrt{11}\text{ rad/s}$, yielding the ultimate period:
\[
  P_u = \frac{2\pi}{\omega_u} = \frac{2\pi}{\sqrt{11}} \approx 1.8944\text{ s}. \qedhere
\]
\end{proof}

Applying the standard Ziegler--Nichols closed-loop PID tuning rules ($K_p = 0.6K_u$, $T_i = 0.5P_u$, and $T_d = 0.125P_u$) yields the analytical baseline gains:
\begin{align}
  K_{p,\text{ZN}} &= 0.6(60) = 36.000, \label{eq:zn_kp} \\
  K_{i,\text{ZN}} &= \frac{K_{p,\text{ZN}}}{0.5 P_u} = \frac{36\sqrt{11}}{\pi} \approx 38.006, \label{eq:zn_ki} \\
  K_{d,\text{ZN}} &= K_{p,\text{ZN}}(0.125 P_u) = \frac{9\pi}{\sqrt{11}} \approx 8.525. \label{eq:zn_kd}
\end{align}

\subsection{Realizable Filtered PID Transfer Function}
An ideal parallel PID controller $C_{\text{ideal}}(s) = K_p + \frac{K_i}{s} + K_d s$ is improper (numerator degree exceeds denominator degree), causing unbounded high-frequency gain ($\lim_{\omega \to \infty} |C_{\text{ideal}}(j\omega)| = \infty$) that amplifies measurement noise. Both MATLAB (\texttt{pid(Kp, Ki, Kd, 0.01)}) and Simulink (\texttt{N = 100}) implement a practical first-order derivative filter with time constant $T_f = 1/N = 0.01\text{ s}$:
\begin{equation}
  C(s) = K_p + \frac{K_i}{s} + \frac{K_d s}{T_f s + 1} = \frac{(K_p T_f + K_d)s^2 + (K_p + K_i T_f)s + K_i}{s(T_f s + 1)} = \frac{N_c(s)}{s(T_f s + 1)},
  \label{eq:pid_filtered}
\end{equation}
where $N_c(0) = K_i$. This renders $C(s)$ biproper and bounds its high-frequency asymptote to $K_p + K_d/T_f = K_p + 100K_d$.

\subsection{Proofs of Zero Steady-State Error and Load Disturbance Rejection}
Consider the two-degree-of-freedom closed-loop response driven by reference input $R(s)$ and plant-input load disturbance $D(s)$:
\begin{equation}
  Y(s) = \underbrace{\frac{C(s)G(s)}{1 + C(s)G(s)}}_{T(s)}\,R(s) + \underbrace{\frac{G(s)}{1 + C(s)G(s)}}_{S_d(s)}\,D(s).
  \label{eq:closed_loop_two_input}
\end{equation}

\begin{proposition}[Asymptotic Step Tracking and Load Disturbance Rejection]
Let $G(s) = \frac{1}{D_p(s)}$ with $D_p(0) = 6$, and let $C(s)$ be given by \eqref{eq:pid_filtered} with $K_i > 0$ such that the closed-loop system is Hurwitz stable. Then:
\begin{enumerate}[label=(\alph*),leftmargin=2em]
  \item For a unit-step reference $R(s) = \frac{1}{s}$ with $D(s) = 0$, the steady-state tracking error is $e_{ss,\text{ref}} = 0$.
  \item For a step load disturbance $D(s) = \frac{A_d}{s}$ ($A_d = -0.3$) with $R(s) = 0$, the steady-state output deviation is $y_{d,ss} = 0$.
\end{enumerate}
\end{proposition}

\begin{proof}
\textbf{(a) Reference Tracking:} The Laplace-domain tracking error is $E(s) = R(s) - Y(s) = \frac{1}{1 + C(s)G(s)}R(s)$. Substituting \eqref{eq:plant} and \eqref{eq:pid_filtered} with $R(s) = \frac{1}{s}$ and applying the Final Value Theorem (FVT), valid because all poles of $sE(s)$ lie in the open left-half plane:
\[
  e_{ss,\text{ref}} = \lim_{t \to \infty} e(t) = \lim_{s \to 0} s E(s) = \lim_{s \to 0} \frac{s(T_f s + 1)D_p(s)}{s(T_f s + 1)D_p(s) + N_c(s)} = \frac{0}{0 + K_i} = 0.
\]

\textbf{(b) Load Disturbance Rejection:} Setting $R(s) = 0$ and $D(s) = \frac{A_d}{s}$ in \eqref{eq:closed_loop_two_input}, the disturbance-induced output $Y_d(s) = S_d(s)D(s)$ becomes:
\[
  Y_d(s) = \frac{G(s)}{1 + C(s)G(s)} \cdot \frac{A_d}{s} = \frac{A_d(T_f s + 1)}{s(T_f s + 1)D_p(s) + N_c(s)}.
\]
Applying the Final Value Theorem yields:
\[
  y_{d,ss} = \lim_{t \to \infty} y_d(t) = \lim_{s \to 0} s Y_d(s) = \lim_{s \to 0} \frac{s A_d(T_f s + 1)}{s(T_f s + 1)D_p(s) + N_c(s)} = \frac{0}{K_i} = 0. \qedhere
\]
\end{proof}

\begin{remark}
The integrator pole at $s = 0$ in $C(s)$ introduces a blocking zero at $s = 0$ in the load sensitivity function $S_d(s) = G(s)/(1 + C(s)G(s))$, guaranteeing complete asymptotic rejection of constant load disturbances while the transient dip magnitude is governed by $|1 + C(j\omega)G(j\omega)|$.
\end{remark}

\subsection{\texorpdfstring{$s$}{s}-Plane Pole Geometry and Frequency-Domain Stability Margins}
For a dominant complex-conjugate closed-loop pole pair $s_{1,2} = -\sigma \pm j\omega_d = -\zeta\omega_n \pm j\omega_n\sqrt{1 - \zeta^2}$, transient performance obeys:
\begin{equation}
  M_p = 100\exp\!\left(-\frac{\pi\zeta}{\sqrt{1 - \zeta^2}}\right), \qquad t_s\,(2\%) \approx \frac{4}{\sigma} = \frac{4}{\zeta\omega_n}, \qquad \zeta = \cos\theta,
\end{equation}
where $\theta$ is the angle between the pole vector and the negative real axis. Shifting dominant poles deeper into the left-half plane (increasing $\sigma$) shortens settling time, while decreasing $\theta$ (increasing damping ratio $\zeta$) suppresses overshoot.

In the frequency domain, let $L(s) = C(s)G(s)$ denote the open-loop transfer function. Relative stability is quantified via:
\begin{itemize}[leftmargin=2em]
  \item \textbf{Phase Margin ($\text{PM}$):} $\text{PM} = 180^\circ + \angle L(j\omega_{cp})$ at the gain crossover frequency $|L(j\omega_{cp})| = 1$.
  \item \textbf{Gain Margin ($\text{GM}_{\text{dB}}$):} $\text{GM}_{\text{dB}} = -20\log_{10}|L(j\omega_{cg})|$ at the phase crossover frequency $\angle L(j\omega_{cg}) = -180^\circ$.
\end{itemize}
Standard textbook design targets $\text{PM} \in [45^\circ, 65^\circ]$ and $\text{GM}_{\text{dB}} > 6\text{ dB}$, whereas Ziegler--Nichols tuning typically yields $\text{PM} \approx 15^\circ\text{--}25^\circ$.

% ==============================================================================
\section{Constrained Multi-Objective PSO Formulation}
% ==============================================================================

\subsection{Search Space and Exterior Hurwitz Stability Barrier}
Each particle $i \in \{1, \dots, 35\}$ encodes a candidate gain vector $\mathbf{x}_i = [K_p, K_i, K_d] \in \Omega$, bounded within $\Omega = [0, 45] \times [0, 45] \times [0, 25]$. In \texttt{src/evaluate\_pid.m}, candidate controllers are evaluated over $t \in [0, 8]\text{ s}$ ($\Delta t = 0.01\text{ s}$) using a composite objective with an exterior stability barrier:
\begin{equation}
  J(\mathbf{x}) = \begin{cases}
    1.0\,\text{ITAE} + 0.6\,M_p + 1.5\,t_s + 25.0\,e_{ss}, & \text{if } \max_m \operatorname{Re}(\lambda_m) < 0, \\[4pt]
    10^6, & \text{otherwise},
  \end{cases}
  \label{eq:cost_function}
\end{equation}
where $\lambda_m$ are the closed-loop eigenvalues of $T(s)$ (\texttt{isstable(T)}) and
\begin{equation}
  \text{ITAE} = \int_0^{T_{\text{opt}}} t\,|1 - y(t)|\,dt
\end{equation}
weights tracking error by elapsed time $t$ so that late-time oscillations are penalized far more severely than the unavoidable initial error at $t = 0$.

\subsection{Swarm Kinematic Update Equations}
Initialized with a fixed deterministic generator (\texttt{rng(42, 'twister')}) across $N_p = 35$ particles and $K_{\max} = 60$ iterations, particle velocities and positions update at iteration $k$ via:
\begin{align}
  \mathbf{v}_i^{k+1} &= w(k)\mathbf{v}_i^k + c_1 \mathbf{r}_1 \odot (\mathbf{p}_{\text{best},i} - \mathbf{x}_i^k) + c_2 \mathbf{r}_2 \odot (\mathbf{g}_{\text{best}} - \mathbf{x}_i^k), \label{eq:pso_vel} \\
  \mathbf{x}_i^{k+1} &= \operatorname{clip}\!\left(\mathbf{x}_i^k + \mathbf{v}_i^{k+1},\; \mathbf{x}_{\min},\; \mathbf{x}_{\max}\right), \label{eq:pso_pos}
\end{align}
where $c_1 = c_2 = 1.5$, $\mathbf{r}_1, \mathbf{r}_2 \sim \mathcal{U}(0,1)^3$, $\odot$ denotes element-wise multiplication, and inertia weight $w(k)$ anneals linearly from $w_{\max} = 0.9$ (global exploration) to $w_{\min} = 0.4$ (local refinement):
\begin{equation}
  w(k) = w_{\max} - \frac{w_{\max} - w_{\min}}{K_{\max}}\,k.
\end{equation}

\subsection{Modular Software and Simulink Validation Architecture}
The repository in \texttt{D:\textbackslash MATLAB\_Projects\textbackslash PID Optimization} separates analytical synthesis from block-diagram verification:
\begin{itemize}[leftmargin=2em]
  \item \texttt{src/evaluate\_pid.m}: Evaluates closed-loop stability, ITAE, overshoot, settling time, and steady-state error.
  \item \texttt{src/pso\_pid\_tuning.m}: Runs PSO optimization, computes Gain and Phase margins via \texttt{margin()}, simulates a step load disturbance ($A_d = -0.3$ at $t = 4\text{ s}$), and exports 300-DPI plots to \texttt{results/}.
  \item \texttt{src/build\_simulink\_model.m}: Programmatically generates \texttt{simulink/pid\_pso\_demo.slx} featuring parallel PSO-PID and Ziegler--Nichols closed loops sharing a unit-step setpoint (\texttt{Setpoint\_Ref}), plant-input disturbance (\texttt{Load\_Disturbance}), and multi-trace \texttt{Comparison\_Scope}.
\end{itemize}

\includeResultFigure{../results/simulink_model.png}{Closed-loop Simulink block diagram (\texttt{pid\_pso\_demo.slx}) comparing PSO-Tuned PID and Ziegler--Nichols PID under unit-step setpoint tracking and plant-input load disturbance ($A_d = -0.3$ at $t = 4\text{ s}$).}{fig:simulink_model}

% ==============================================================================
\section{Textbook Curriculum Alignment and Method Comparison}
% ==============================================================================

\subsection{Mapping to Canonical Control Systems Literature}
Table~\ref{tab:textbook_mapping} maps each implementation module in this project to its theoretical foundation in standard undergraduate and graduate control engineering textbooks.

\begin{table}[H]
  \centering
  \caption{Alignment of Project Components with Standard Control Systems and Optimization Textbooks}
  \label{tab:textbook_mapping}
  \small
  \begin{tabularx}{\textwidth}{>{\raggedright\arraybackslash}p{0.23\textwidth} >{\raggedright\arraybackslash}p{0.23\textwidth} >{\raggedright\arraybackslash}p{0.27\textwidth} >{\raggedright\arraybackslash}X}
    \toprule
    \textbf{Project Module} & \textbf{Textbook Concept} & \textbf{Mathematical Basis} & \textbf{Canonical Reference} \\
    \midrule
    \texttt{G = tf(1, [1 6 11 6])} & Higher-Order LTI Modeling & Continuous transfer function with real negative poles & Ogata (Ch.\ 2, 5); Nise (Ch.\ 4) \\
    \addlinespace[3pt]
    \texttt{Ku = 60}, \texttt{Pu = 2*pi/sqrt(11)} & Routh--Hurwitz \& Ultimate Sensitivity & Marginal stability on $j\omega$-axis; quarter-decay-ratio rules & Ogata (Ch.\ 6, 8) \\
    \addlinespace[3pt]
    \texttt{pid(Kp, Ki, Kd, 0.01)} & Realizable Filtered PID & First-order derivative filter $T_f = 1/N$ limiting noise gain & \AA str\"om \& H\"agglund (Ch.\ 3) \\
    \addlinespace[3pt]
    \texttt{stepinfo(y, t, 1)} & Time-Domain Specifications & Percent overshoot $M_p$, $2\%$ settling time $t_s$, steady-state error & Dorf \& Bishop (Ch.\ 5); Nise (Ch.\ 4) \\
    \addlinespace[3pt]
    \texttt{trapz(t, t.*abs(1-y))} & Optimal Performance Indices & Integral of Time-weighted Absolute Error ($\text{ITAE}$) & Dorf \& Bishop (Ch.\ 5); Ogata (Ch.\ 8) \\
    \addlinespace[3pt]
    \texttt{margin(C * G)} & Frequency Stability Margins & Gain Margin ($\text{GM}_{\text{dB}}$) and Phase Margin ($\text{PM}^\circ$) & Ogata (Ch.\ 9); Nise (Ch.\ 10) \\
    \addlinespace[3pt]
    \texttt{pzplot(T\_zn, T\_pso)} & Root Locus \& $s$-Plane Geometry & Left-half-plane pole migration and damping ratio $\zeta = \cos\theta$ & Nise (Ch.\ 8); Dorf \& Bishop (Ch.\ 7) \\
    \addlinespace[3pt]
    \texttt{Load\_Disturbance} ($t=4\text{ s}$) & Sensitivity \& Disturbance Rejection & Load sensitivity $S_d(s) = \frac{G(s)}{1+C(s)G(s)}$ and FVT offset rejection & \AA str\"om \& Murray (Ch.\ 11) \\
    \addlinespace[3pt]
    \texttt{pso\_pid\_tuning.m} & Swarm Intelligence Optimization & Velocity--position updates with linearly decreasing inertia $w(k)$ & Kennedy \& Eberhart (1995); Gaing (2004) \\
    \bottomrule
  \end{tabularx}
\end{table}

\subsection{Comparative Evaluation of PID Tuning Methodologies}
Table~\ref{tab:method_comparison} contrasts classical Ziegler--Nichols rules and analytical pole placement against the multi-objective PSO approach implemented in this project.

\begin{table}[H]
  \centering
  \caption{Theoretical Comparison of Classical, Analytical, and Swarm-Based PID Tuning Methods}
  \label{tab:method_comparison}
  \small
  \begin{tabularx}{\textwidth}{>{\raggedright\arraybackslash}p{0.20\textwidth} >{\raggedright\arraybackslash}X >{\raggedright\arraybackslash}X >{\raggedright\arraybackslash}X}
    \toprule
    \textbf{Criterion} & \textbf{Ziegler--Nichols (Closed-Loop)} & \textbf{Analytical Pole Placement} & \textbf{PSO-Tuned PID (This Work)} \\
    \midrule
    \textbf{Design Basis} & Empirical formulas from ultimate gain $K_u$ and period $P_u$ & Algebraic matching of closed-loop characteristic polynomial roots & Stochastic global search over a weighted time-domain cost surface \\
    \addlinespace[3pt]
    \textbf{Plant Knowledge} & Requires marginal oscillation test ($K_u, P_u$) & Requires exact symbolic plant order and zero-pole structure & Requires only black-box closed-loop step simulation \\
    \addlinespace[3pt]
    \textbf{Overshoot ($M_p$)} & High ($40\%\text{--}60\%$) due to quarter-decay-ratio calibration & Exact for 2nd-order plants; degraded by unmodeled dynamics & Directly minimized via explicit $M_p$ and $\text{ITAE}$ objective weights \\
    \addlinespace[3pt]
    \textbf{Stability Margins} & Narrow ($\text{PM} \approx 15^\circ\text{--}25^\circ$), prone to oscillatory ringing & Indirectly governed by chosen damping ratio $\zeta$ & Wide ($\text{PM} \in [45^\circ, 65^\circ]$, $\text{GM} > 6\text{ dB}$) via stability barrier \\
    \addlinespace[3pt]
    \textbf{Disturbance Recovery} & Fast initial response followed by prolonged oscillation & Trade-off between loop bandwidth and sensor noise amplification & Small transient drop with rapid, non-oscillatory return to setpoint \\
    \bottomrule
  \end{tabularx}
\end{table}

% ==============================================================================
\section{Experimental Results and Validation}
% ==============================================================================

\subsection{Quantitative Time-Domain and Frequency-Domain Benchmarks}
Tables~\ref{tab:time_results} and~\ref{tab:freq_results} organize the comparative metrics produced by \texttt{src/pso\_pid\_tuning.m}. Analytical Ziegler--Nichols gains from \eqref{eq:zn_kp}--\eqref{eq:zn_kd} are pre-populated; populate the remaining simulation entries directly from the MATLAB Command Window printout after executing \texttt{pso\_pid\_tuning.m}.

\begin{table}[H]
  \centering
  \caption{Synthesized PID Gains, Step-Response Metrics, and Load Disturbance Drop ($A_d = -0.3$ at $t = 4\text{ s}$)}
  \label{tab:time_results}
  \small
  \begin{tabular}{lccccccc}
    \toprule
    \textbf{Tuning Method} & $K_p$ & $K_i$ & $K_d$ & \textbf{Overshoot} & \textbf{Settling} & \textbf{SS Error} & \textbf{Dist.\ Drop} \\
    \midrule
    Ziegler--Nichols PID & $36.000$ & $38.006$ & $8.525$ & $44.94\%$ & $4.94$ & $0.0018$ & $0.0338$ \\
    PSO-Tuned PID        & $18.002$ & $10.627$ & $7.880$ & $0.04\%$ & $1.20$ & $0.0001$ & $0.0110$ \\
    \bottomrule
  \end{tabular}
\end{table}

\begin{table}[H]
  \centering
  \caption{Open-Loop Frequency-Domain Stability Margins Computed via \texttt{margin(C * G)}}
  \label{tab:freq_results}
  \small
  \begin{tabular}{lcccc}
    \toprule
    \textbf{Tuning Method} & \textbf{Gain Margin (dB)} & \textbf{Phase Margin (deg)} & $\omega_{cg}$ \textbf{(rad/s)} & $\omega_{cp}$ \textbf{(rad/s)} \\
    \midrule
    Ziegler--Nichols PID & $26.73$ & $28.98$ & $13.72$ & $2.64$ \\
    PSO-Tuned PID        & $33.69$ & $69.63$ & $19.42$ & $1.86$ \\
    \bottomrule
  \end{tabular}
\end{table}

\subsection{Visual Step Tracking, Pole-Zero Migration, and Swarm Convergence}
Figures~\ref{fig:step_disturbance}, \ref{fig:pole_zero}, and \ref{fig:convergence} display the exported 300-DPI validation plots from \texttt{results/}:

\includeResultFigure{../results/step_disturbance_comparison.png}{Closed-loop unit-step tracking ($t = 0\text{ s}$) and plant-input load disturbance rejection ($A_d = -0.3$ at $t = 4\text{ s}$) comparing Ziegler--Nichols PID and PSO-Tuned PID.}{fig:step_disturbance}

\includeResultFigure{../results/simulink_scope_output.png}{Simulink \texttt{Comparison\_Scope} time-domain waveform validating unit-step setpoint tracking ($t = 0\text{ s}$) and plant-input load disturbance recovery ($A_d = -0.3$ at $t = 4\text{ s}$) simulated with the \texttt{ode45} solver.}{fig:simulink_scope}

\includeResultFigure{../results/pole_zero_map.png}{Closed-loop $s$-plane pole-zero map comparing left-half-plane damping migration between Ziegler--Nichols PID and PSO-Tuned PID.}{fig:pole_zero}

\includeResultFigure{../results/convergence_curve.png}{Global best multi-objective fitness cost $J(\mathbf{g}_{\text{best}})$ across 60 PSO iterations ($N_p = 35$ particles).}{fig:convergence}

% ==============================================================================
\section{Engineering Trade-Offs and Conclusions}
% ==============================================================================

Four core control-engineering principles govern the observed performance gains of the PSO-tuned controller:
\begin{enumerate}[leftmargin=2em]
  \item \textbf{Damping Ratio vs.\ Quarter-Decay Ratio:} Ziegler--Nichols tuning enforces a quarter-amplitude decay ratio ($e^{-2\pi\zeta/\sqrt{1-\zeta^2}} = 0.25 \implies \zeta \approx 0.215$), which inherently produces $\approx 50\%$ overshoot and a low phase margin. By penalizing $M_p$ and $t_s$ directly in \eqref{eq:cost_function}, PSO shifts dominant closed-loop poles to a higher damping ratio $\zeta$, eliminating sustained oscillations.
  \item \textbf{ITAE vs.\ ISE Performance Indices:} Because initial step error at $t = 0$ is fixed at $e(0) = 1$, Integral of Squared Error ($\text{ISE} = \int e^2\,dt$) favors high-gain aggressive controllers that overshoot. Weighting error by elapsed time ($\text{ITAE} = \int t|e(t)|\,dt$) prioritizes rapid settling and oscillation suppression.
  \item \textbf{Hurwitz Stability Barrier:} Checking \texttt{isstable(T)} inside \texttt{src/evaluate\_pid.m} maps the Routh--Hurwitz stability boundary into a hard penalty ($J = 10^6$), preventing particle trajectories from crossing into the right-half $s$-plane.
  \item \textbf{Practical Derivative Filtering:} Incorporating the first-order pole at $s = -1/T_f = -100\text{ rad/s}$ preserves derivative phase lead near the gain crossover frequency while rolling off high-frequency sensor noise amplification.
\end{enumerate}

Combining analytical Routh--Hurwitz stability bounds, Final Value Theorem proofs, and multi-objective swarm optimization yields a reproducible, textbook-grounded methodology for automated PID synthesis in MATLAB and Simulink.

% ==============================================================================
\begin{thebibliography}{9}
  \bibitem{ogata2010}
  K.~Ogata, \textit{Modern Control Engineering}, 5th~ed. Upper Saddle River, NJ, USA: Prentice Hall, 2010.

  \bibitem{nise2019}
  N.~S. Nise, \textit{Control Systems Engineering}, 8th~ed. Hoboken, NJ, USA: Wiley, 2019.

  \bibitem{dorf2017}
  R.~C. Dorf and R.~H. Bishop, \textit{Modern Control Systems}, 13th~ed. Hoboken, NJ, USA: Pearson, 2017.

  \bibitem{astrom1995}
  K.~J. \AA str\"om and T.~H\"agglund, \textit{PID Controllers: Theory, Design, and Tuning}, 2nd~ed. Research Triangle Park, NC, USA: Instrument Society of America, 1995.

  \bibitem{astrom2008}
  K.~J. \AA str\"om and R.~M. Murray, \textit{Feedback Systems: An Introduction for Scientists and Engineers}. Princeton, NJ, USA: Princeton University Press, 2008.

  \bibitem{kennedy1995}
  J.~Kennedy and R.~Eberhart, ``Particle swarm optimization,'' in \textit{Proc. IEEE Int. Conf. Neural Networks (ICNN)}, Perth, WA, Australia, Nov. 1995, vol.~4, pp.~1942--1948.

  \bibitem{gaing2004}
  Z.-L. Gaing, ``A particle swarm optimization approach for optimum design of PID controller in AVR system,'' \textit{IEEE Trans. Energy Conversion}, vol.~19, no.~2, pp.~384--391, Jun. 2004.
\end{thebibliography}

\end{document}